\section{Purpose, conventions, and main conclusions}

\subsection{What is being bounded?}

For an irrational real number \(\xi\), its irrationality exponent is
\[
\mu(\xi)=\sup\left\{u>0:
0<\left|\xi-\frac pq\right|<q^{-u}
\text{ for infinitely many reduced }p/q,\ q>0\right\}.
\]
Thus \(\mu(\xi)\le E\) means that for every \(\varepsilon>0\) there are
only finitely many reduced fractions with error below
\(q^{-E-\varepsilon}\). Equivalently, after changing a positive
constant, one has
\begin{equation}\label{eq:measuremeaning}
\left|\xi-\frac pq\right|\ge c_\varepsilon q^{-E-\varepsilon}
\end{equation}
for all \(p,q\) with \(q>0\). This does not assert the same inequality
with exponent exactly \(E\).

Dirichlet's theorem gives \(\mu(\xi)\ge2\). Obtaining an explicit upper
bound for a particular transcendental number is a separate problem.
The finiteness of the exponents of fixed powers of \(\pi\) is already
accessible through classical quantitative transcendence theory.
Our objective is to spell out one concrete construction and track
its constants, not to claim a first finiteness theorem.

All logarithms in estimates are natural logarithms. We use
\[
\lambda=\log\ii=\frac{\ii\pi}{2},\qquad \Theta=\frac\pi2.
\]
The symbols \(o(n)\) and \(O(\log n)\) refer to \(n\to\infty\)
with the complete interval configuration and target \(k\) fixed.
The auxiliary degree \(m\) may be large, but is fixed in this limit.

\subsection{Why odd powers are natural targets}

The Dirichlet beta function is
\[
\beta(s)=\sum_{\ell=0}^{\infty}\frac{(-1)^\ell}{(2\ell+1)^s}.
\]
For \(r\ge0\),
\[
\beta(2r+1)=
\frac{(-1)^r E_{2r}}{4^{r+1}(2r)!}\pi^{2r+1},
\]
where \(\operatorname{sech}z=\sum_{j\ge0}E_jz^j/j!\) defines the
Euler numbers. In particular,
\[
\beta(1)=\frac\pi4,\qquad
\beta(3)=\frac{\pi^3}{32},\qquad
\beta(5)=\frac{5\pi^5}{1536}.
\]
Multiplication by a nonzero rational number preserves the irrationality
exponent. For a reduced multiplier \(a/b\), reducing \(ap/(bq)\)
cancels a gcd dividing \(|ab|\), so its denominator remains comparable
to \(q\) by constants independent of \(p,q\).

This observation does \emph{not} deduce an upper bound for \(\pi^3\)
from a rational-approximation bound for \(\pi\). Taking a cube root
changes the approximation field. The resultant argument below accounts
for that change.

\subsection{A preview of the construction}

There are three components:
\begin{enumerate}
\item Build integer polynomials \(A_{h,j,n}(x)\) such that
\[
R_{h,n}(x)=\sum_{j=0}^{m}A_{h,j,n}(x)(\log x)^j
\]
has a high-order zero at \(x=1\).
\item Divide every coefficient by a common integer \(\Pi_n\).
Its exponential rate \(\gamma=\lim n^{-1}\log\Pi_n\) matters;
a large constant depending only on \(m\) does not change this rate.
\item Evaluate at \(x=\ii\), compare \(\lambda\) with a \(k\)-th root
of a rational approximation to \(\pi^k\), and use a nonzero resultant.
\end{enumerate}
Nested intervals change both analytic bounds and arithmetic divisors.
Their effects must be assessed together.

The classical logarithmic construction appears in
Mignotte~\cite{Mignotte}. Weighted logarithmic Pad\'e systems and
normality are discussed in Kawashima~\cite{Kawashima}.
Here the determinant and estimates for our two-sided nested system
are proved directly.

\subsection{General parameter theorem}

Choose integers
\[
L\ge3,\quad 0=t_0<t_1<\cdots<t_{J-1}<L/2,\quad \nu_j\ge1,
\]
and put
\[
M=\sum_{j=0}^{J-1}\nu_j,\qquad m=M-1\ge1,\qquad
c_j=1-\frac{2t_j}{L}.
\]
An interval of relative length \(c_j\), centred at \(1/2\), is used
\(\nu_j\) times. Let \(\gamma\) be the explicit prime-divisor rate
defined in Section~\ref{sec:arithmetic}. Let \(Y>0\) solve
\begin{equation}\label{eq:rootpreview}
2\sum_j\nu_j\arctan\frac{c_j}{2Y}=\Theta.
\end{equation}
Define
\begin{align}
F_c(Y)&=\frac c2\log\frac{1+(2Y/c)^2}{4}
             +2Y\arctan\frac{c}{2Y},\label{eq:Fdef}\\
D_0&=\sum_j\nu_jF_{c_j}(Y)-\Theta Y-m,\label{eq:D0def}\\
G_0&=m+\left(\sum_j\nu_jc_j\right)\log2,\label{eq:G0def}\\
D&=D_0+\gamma,\qquad G=G_0-\gamma.\label{eq:DGdef}
\end{align}

\begin{theorem}[Odd-power bound]\label{thm:main}
Let \(k\ge3\) be odd. If \(G\ge0\) and
\(\Delta_k=D-(k-1)G>0\), then
\begin{equation}\label{eq:mainbound}
\boxed{\displaystyle
\mu(\pi^k)\le E_k:=
\frac{m(D+G)}{D-(k-1)G}
=\frac{m(D_0+G_0)}
 {D_0-(k-1)G_0+k\gamma}.}
\end{equation}
\end{theorem}

The proof also works for every integer \(k\ge2\). We emphasize odd
powers because they are the motivating targets; we do not compare
these bounds with specialized results for even powers.

\begin{corollary}[A twelve-layer family]\label{cor:numerical}
Use \(L=69\), \(t_j=j\) for \(0\le j\le11\), and
\[
\nu_0=6s,\qquad \nu_1=\cdots=\nu_{11}=s,\qquad m=17s-1.
\]
The following choices satisfy the hypotheses:
\begin{center}
\begin{tabular}{rrrr}
\toprule
\(k\)&\(s\)&\(m\)&Certified strict upper bound\\
\midrule
3&11&186&\(1203.973\)\\
5&259&\(4\,402\)&\(42\,461\)\\
7&\(5\,871\)&\(99\,806\)&\(1\,272\,103\)\\
9&\(131\,354\)&\(2\,233\,017\)&\(35\,398\,937\)\\
\bottomrule
\end{tabular}
\end{center}
For every other fixed odd \(k\), some sufficiently large integer \(s\)
also satisfies the hypotheses.
\end{corollary}

The last assertion is an existence statement within an explicit
family, not a uniform small bound as \(k\to\infty\).
Section~\ref{sec:parameters} explains the growth.

\subsection{A guide to the evidence}

The algebraic and analytic lemmas prove Theorem~\ref{thm:main}.
Ball-arithmetic certificates prove the decimal comparisons once the
lemmas hold. Exact finite examples check identities and implementation,
but do not prove asymptotic rates. Floating-point searches select
parameters; they neither certify rounding nor prove global optimality.
The appendices retain unsuccessful experiments because they explain
the arithmetic difficulties that remain.
